% ===================================================================
%  اللوحة نمرة ٧ — القرية في الشتا. البيت الريفي في الربيع.
%
%  Traduction, faite en 2026, en arabe egyptien ecrit, sur l'ido de
%  texto/io. Regles de la colonne : voir 10-lawha-01.tex. Deux
%  scenes, deux numerotations : les cles portent c1 et c2.
%
%  LA CAMPAGNE EST L'ENDROIT OU L'EGYPTIEN A LE PLUS DE MOTS A LUI,
%  PARCE QUE C'EST LA QU'IL A LE PLUS DE CHOSES :
%
%    غيط (3, 4)    le champ        MSA : حقل
%    جرن (33)      la grange, l'aire a battre
%    زريبة (48)    la bergerie
%    برج الحمام (42) le pigeonnier
%    خُمّ (44)      le poulailler
%    عزبة (36, 60) la ferme
%    عشّة (38)      la niche, la hutte
%    زبل (81)      le fumier
%    بير (29)      le puits        MSA : بئر
%    ترعة (23)     le ruisseau, le canal
%
%  DEUX DE CES MOTS TOMBENT SUR LA PLANCHE MIEUX QU'UNE TRADUCTION.
%  La roue du moulin (17) est ساقية, qui est le nom de la roue a eau
%  egyptienne : meme objet, meme fonction, et le mot existait avant
%  qu'on ait a le chercher. Et le pigeonnier (42) est برج الحمام, la
%  tour a pigeons, qui est un point de repere de tout village
%  egyptien — la gravure francaise en montre un, et il n'a pas fallu
%  le nommer par periphrase.
%
%  L'ECOLE DU VILLAGE APPORTE SON MOT, LUI AUSSI VENU DU TURC :
%
%    أبلة (13)   la maitresse d'ecole — de abla, la grande sœur
%
%  et les metiers continuent en -جي : الجزمجي (7) le cordonnier,
%  après العربجي et الكهربجي du tableau 5.
%
%  LE RESTE DU VILLAGE :
%
%    لوكاندة (2)  l'auberge   de l'italien locanda
%    يافطة (3)    l'enseigne
%    مدافن (20)   le cimetiere  MSA : مقبرة
%    خفير (27)    le garde champetre
%    أبو فروة (42) le marron    MSA : كستناء
%    موس          le rasoir
%
%  UN SEUL BLOC A DEMANDE UNE REPRISE : la glace du ruisseau
%  (c1-10-1), ou l'ido nomme la glace avant l'eau qui la porte.
% ===================================================================

%%K t07-noto-1 noto
\VUnotes{143.2mm}{%
(*) بالنسبة لتفاصيل الأكل، بُصّ اللوحتين نمرة ٦ ونمرة ١٣.%
}

%%K t07-apar-1 apar
\VUsaut{4.0mm}
\VUcentre{13.2pt}{90}{السلسلة التانية}
\VUsaut{3.0mm}
\VUfilet{16mm}
\VUsaut{5.6mm}
\VUtitre{15.0pt}{60}{اللوحة نمرة ٧}
\VUsaut{3.0mm}

%%K t07-tit-1 sub
\VUcentre{12.0pt}{0}{\textit{المشهد الأول.}}
\VUsaut{3.0mm}

%%K t07-tit-2 sub
\VUpk{10.2pt}{40}{القرية في الشتا.}
\VUsaut{4.0mm}

%%K t07-c1-01-1 p
1. --- في أجازة رأس السنة، رحت مع عمّي لمدينة نوفيل، قرية صغيرة كان عنده فيها
شوية شغل تجاري.

%%K t07-c1-02-1 p
2. --- ركبنا \VUgras{الدليجانس} \textsuperscript{(11)} اللي بيشتغل بين المدينة
والقرية دي؛ بس علشان الدنيا كانت ساقعة أوي، السفرية دي في عربية مش مريحة أوي
ما كانتش لذيذة قوي.

%%K t07-c1-03-1 p
3. --- علشان كده فرحنا أوي لما شفنا \VUgras{العربجي} بيوقّف عربيته في الميدان،
قدام \VUgras{لوكاندة} \textsuperscript{(2)} \textit{الدبّ الأبيض}. و\VUgras{صاحب
اللوكاندة} \textsuperscript{(4)} كان مستنّينا على عتبة بيته، وبيشرب \VUgras{غليونه}
بهدوء.

%%K t07-c1-03-2 p
عزمنا ندخل، وأنا ما استنّيتش حد يعزمني تاني علشان أقعد قدام نار الحطب اللي
كانت بتفرقع في الموقد. ساعتها ضحكت أوي على \VUgras{الهوا الشمالي} القارص اللي
بيخلّي \VUgras{اليافطة} \textsuperscript{(3)} القديمة تزنّ، وبيشيل بدوّامات سودا
\VUgras{الدخان} \textsuperscript{(1)} اللي طالع من المدخنة.

%%K t07-c1-04-1 p
4. --- وكانت ساعة الغدا. من غير ما نستنّى أكتر، أكلنا كويس الأكل البسيط بس
اللذيذ اللي قدّموه لنا (*).

%%K t07-tit-3 sub
\VUpk{10.2pt}{40}{كام زيارة.}
\VUsaut{3.0mm}

%%K t07-c1-05-1 p
5. --- بعد الغدا، رحت مع عمّي، اللي بيشتغل في التأمين، عند زباينه. وأول واحد
زرناه كان \VUgras{الحدّاد} \textsuperscript{(5)}، اللي ساكن جنب اللوكاندة. كان مشغول
بتنعيل حصان. وأنا اتبهرت من قوّة صاحبه اللي كان ماسك بثبات على مريلته الجلد
\VUgras{حافر} الحصان، والحدّاد بيثبّت \VUgras{النعل} \textsuperscript{(6)} بضربات
شاكوش قوية.

%%K t07-c1-06-1 p
6. --- وبعدين رحنا عند \VUgras{الجزمجي} \textsuperscript{(7)} بتاع القرية، يعقوب
العجوز. ومع إن يافطته \VUgras{بوت} حلو أوي، كان راضي إنه يركّب نعل جديد لجوز
جزم قديم مخرّم.

%%K t07-c1-06-2 p
والعجوز قال لنا وهو بيضحك: «في المدينة، كنت اسمي (صانع أحذية) وما كانش عندي
زباين. هنا، اسمي (جزمجي) والشغل كتير.»

%%K t07-c1-07-1 p
7. --- وحكى لنا عن جاره \VUgras{الحلّاق} \textsuperscript{(8)}، اللي زباينه مش
راضيين عنه. \VUgras{مواسه} على طول متكسّرة السنّ و\VUgras{مقصّاته} ولا مرة
بتقصّ كويس.

بس علشان هو الحلّاق الوحيد في القرية، الناس مضطرّة طبعًا تحلق عنده وتقصّ شعرها
عنده. وكمان هو راجل بخيل وتقيل الدم.

%%K t07-c1-08-1 p
8. --- ولحسن الحظ، باقي \VUgras{الجيران} مش زيه. مثلًا \VUgras{صانع العربيات}
\textsuperscript{(9)} راجل طيب، وشغّيل، وبيحبّ يساعد. تصليح عربية عنده متعة حقيقية.
وشغله دايمًا بيبقى متقن.

%%K t07-c1-09-1 p
9. --- ويعقوب العجوز كمان ما اشتكاش من \VUgras{السرّاج} \textsuperscript{(10)}، اللي
بيساعده أحيانًا في خياطة \VUgras{السروج} لما الشغل يبقى مستعجل.

%%K t07-c1-09-2 p
وعمّي سأله عن عيلته: «كلهم كويسين. حفيدتي في \VUgras{المدرسة} \textsuperscript{(12)}.
و\VUgras{الأبلة} \textsuperscript{(13)} بتاعتها راضية عنها أوي، وبتقول إنها
\VUgras{تلميذة} \textsuperscript{(14)} كويسة. وهي مش زي ابني الشقي؛ ده ما بيفكّرش غير
في اللعب. و\VUgras{المدرِّس} \textsuperscript{(15)} مش قادر يعلّمه أي حاجة.»

%%K t07-tit-4 sub
\VUsaut{3.0mm}
\VUpk{10.2pt}{40}{كلام القرية.}
\VUsaut{3.0mm}

%%K t07-c1-10-1 p
10. --- وبعدين حكى لنا العجوز أخبار القرية. هيبنوا مدرسة جديدة، علشان اللي
موجودة في مجلس القرية بقت صغيرة أوي. وده سهل، علشان هيكفي إنهم يشتروا جزء من
جنينة الطحّان. وده هيبقى مكسب كويس للراجل المسكين. وعلشان البرد، \VUgras{حجر}
\VUgras{الطاحونة} \textsuperscript{(16)} ما بقاش يقدر يطحن القمح، علشان
\VUgras{الساقية} \textsuperscript{(17)} وقفت من \VUgras{التلج} \textsuperscript{(25)} اللي
جمّد \VUgras{الترعة} \textsuperscript{(23)}.

%%K t07-c1-11-1 p
11. --- «وبمناسبة البناء، شفت \VUgras{كنيستنا} \textsuperscript{(18)} الجديدة؟ شكلها
حلو زي الصورة تحت غطاها التلج. و\VUgras{الغربان} \textsuperscript{(19)}، اللي ما
بقتش عارفة برجها القديم، واقفة زعلانة على الشجرة الكبيرة بتاعة \VUgras{المدافن}
\textsuperscript{(20)}.»

%%K t07-c1-12-1 p
12. --- وفكرة المدافن دي فكّرت الجزمجي بالناس اللي عمّي كان يعرفهم واللي ماتوا
السنة اللي فاتت. «يا سلام، كام \VUgras{قبر} \textsuperscript{(21)} يا سيدي اتزادوا
على التانيين، تحت \VUgras{شجرة السرو} \textsuperscript{(22)}! الموت حصد كاتب العدل
العجوز بتاعنا. وأهل القرية كلهم حضروا \VUgras{جنازته}.»

%%K t07-c1-13-1 p
13. --- «أهو اللي خد مكانه بيتزحلق على الترعة. جه لي حالًا يصلّح رباط
\VUgras{الزلّاجة} \textsuperscript{(26)} بتاعته اللي كان اتقطع.»

%%K t07-c1-14-1 p
14. --- «وأهو كمان على حافّة الترعة، \VUgras{الخفير} \textsuperscript{(27)} الجديد.
ده كان عسكري قديم وبيخوّف عيال القرية الأشقيا. وهما بيجروا على طول أول ما
يشوفوا \VUgras{لفّاته} \textsuperscript{(29)} و\VUgras{كاسكيته} \textsuperscript{(28)}.»

%%K t07-c1-15-1 p
15. --- «آه! أهو يوسف العجوز، سايق \VUgras{عربية صغيرة فيها حزم حطب}
\textsuperscript{(30)} للفرّان. --- صحيح، قال عمّي، أنا عارف عدّة \VUgras{البغال}
\textsuperscript{(31)} بتاعته. أنا عايز أكلّمه، هستغلّ الفرصة. لحدّ ما نشوفك يا عمّ
يعقوب.»

%%K t07-c1-16-1 p
16. --- وأول ما خرجنا، عيال القرية سابوا المدرسة وجريوا على \VUgras{المزلقان}
\textsuperscript{(33)}، وهما مش خايفين إنهم يقعوا ويكسروا حاجة في \VUgras{الوقعة}
\textsuperscript{(34)}. وواحد فيهم خد \VUgras{الزحّافة} \textsuperscript{(32)} بتاعته من
عند صانع العربيات، وقعّد فيها واحد من زمايله، وجري بيها.

%%K t07-c1-17-1 p
17. --- والباقيين جريوا على \VUgras{كومة التلج} \textsuperscript{(37)} اللي جنب
\VUgras{الكوبري} \textsuperscript{(40)}، وابتدوا يقصفوا \VUgras{راجل التلج}
\textsuperscript{(39)} اللي عملوه إمبارح واللي حطّوا له برنيطة وغليون وشنطة مدرسة
على كتفه.

%%K t07-c1-18-1 p
18. --- وهما مش واخدين بالهم أوي من الهوا اللي بيجمّد ولا من البرد. وأغلبهم
حتى ما عندوش \VUgras{كوفية} \textsuperscript{(35)}؛ وعلشان يدفّوا بعد معركة
\VUgras{كور التلج} \textsuperscript{(38)}، بتكفيهم حفنة \VUgras{أبو فروة}
\textsuperscript{(42)} سخنة أوي، مشتريينها من العجوزة جاكوبينا \VUgras{البايعة}،
اللي بترتعش من البرد تحت \VUgras{شالها} \textsuperscript{(43)} الرقيّق أوي.

%%K t07-tit-5 sub
\VUsaut{3.0mm}
\VUcentre{12.0pt}{0}{\textit{المشهد التاني.}}
\VUsaut{3.0mm}

%%K t07-tit-6 sub
\VUpk{10.2pt}{40}{البيت الريفي في الربيع.}
\VUsaut{4.0mm}

%%K t07-c2-01-1 p
1. --- بالرغم من كل متع الشتا، إحنا بنرحّب برجوع \VUgras{الربيع} بفرحة عميقة.

%%K t07-c2-01-2 p
وما أحلى المنظر اللي بيعمله حوش العزبة، بالليل، في شهر مايو!

%%K t07-c2-02-1 p
2. --- وعند الأفق، \VUgras{الشمس} \textsuperscript{(1)} بتغيب بالراحة، وبتغرّق
\VUgras{الغيطان} \textsuperscript{(3)} في \VUgras{أشعّتها} \textsuperscript{(2)} البنفسجية.
و\VUgras{الحرّاث} \textsuperscript{(6)} النشيط مستعجل علشان يحرث \VUgras{غيطه}
\textsuperscript{(4)}.

%%K t07-c2-02-2 p
وبـ\VUgras{محراثه} \textsuperscript{(5)} المربوط في \VUgras{حصان} \textsuperscript{(7)} قوي،
بيرسم \VUgras{التلم} \textsuperscript{(8)} اللي فيه \VUgras{البذّار} \textsuperscript{(9)}
هيرمي بملء إيده \VUgras{البذور} اللي هتطلّع السنابل الدهبية.

%%K t07-c2-03-1 p
3. --- و\VUgras{الحشّاشين} \textsuperscript{(11)} بمناجلهم حشّوا حشيش المرعى،
واللي بيقلّبوا الحشيش نشّفوا \VUgras{الدريس} \textsuperscript{(10)} وهما بيقلّبوه
بـ\VUgras{المذاري} \textsuperscript{(12)} والأمشاط، وأهم راجعين.

%%K t07-c2-03-2 p
وفي ناس ماشية قدام العربية التقيلة اللي بتجرّها التيران؛ شايلين عدّتهم على
كتافهم. وناس تانية، أكسل شوية، قعدوا بمزاجهم على الدريس اللي ريحته حلوة.

%%K t07-c2-04-1 p
4. --- وهما بيعدّوا بيسلّموا سلام ودّي على \VUgras{الجنايني} \textsuperscript{(16)}
اللي مشغول في \VUgras{جنينة الفاكهة} \textsuperscript{(13)}، بيقلّم \VUgras{شجر
الفاكهة} ويطعّم \VUgras{شجر الكمّترى} \textsuperscript{(14)}. و\VUgras{شجر البرقوق}
\textsuperscript{(15)} بدأ يزهّر، وفي \VUgras{جنينة الخضار} \textsuperscript{(17)}
المسمّدة كويس، الخضار والكرنب والخس كلهم كويسين.

%%K t07-c2-04-2 p
وعلى الحيطة الواطية، \VUgras{طاووس} \textsuperscript{(18)} حلو بيفرد ديله، علشان
الناس تتفرّج على ريشه الحلو أوي.

%%K t07-tit-7 sub
\VUpk{10.2pt}{40}{حوش العزبة.}
\VUsaut{3.0mm}

%%K t07-c2-05-1 p
5. --- أبواب \VUgras{الإسطبل} \textsuperscript{(19)} مفتوحة، وباين المعلف
و\VUgras{الفرش} \textsuperscript{(23)} بتاع \VUgras{الفرس} \textsuperscript{(21)} اللي
\VUgras{السايس} \textsuperscript{(20)} حالًا فكّها. و\VUgras{المهر} \textsuperscript{(22)}
المضحك أوي برجليه الطويلة، بيمشي ورا أمه وهو بينطّ.

%%K t07-c2-06-1 p
6. --- وجنب \VUgras{البير} \textsuperscript{(29)}، \VUgras{البقرات}
\textsuperscript{(25)} رايحة تشرب من \VUgras{الحوض الخشب} \textsuperscript{(26)} اللي
بيستعملوه مشرب. و\VUgras{العجل} \textsuperscript{(24)} مستغلّ الفرصة علشان يرضع،
و\VUgras{التور} \textsuperscript{(27)} وهو شامم ريحة المرعى الحلوة، بيخوّر من الفرحة
ويرفع راسه بفخر بـ\VUgras{قرونه} \textsuperscript{(28)} القوية.

%%K t07-c2-07-1 p
7. --- والكل شغّال. \VUgras{الشغّالة} \textsuperscript{(32)} ماسكة بإيدها
\VUgras{حبل} \textsuperscript{(31)} البير. بتطلّع مية بـ\VUgras{جردل}
\textsuperscript{(30)} علشان تسقي المواشي. وجنب \VUgras{الجرن} \textsuperscript{(33)}،
راجل بيمشّط التبن، وقدام باب \VUgras{البيت} \textsuperscript{(34)} \VUgras{غزّالة}
\textsuperscript{(35)} عجوزة، بدولابها ومغزلها، بتغزل \VUgras{الكتّان} أو
\VUgras{القنّب} زي زمان.

%%K t07-tit-8 sub
\VUsaut{3.0mm}
\VUpk{10.2pt}{40}{صاحب العزبة.}
\VUsaut{3.0mm}

%%K t07-c2-08-1 p
8. --- \VUgras{صاحب العزبة} \textsuperscript{(36)} بيدير الشغل ده كله وبيدّي تعليمات
لعمّاله. وبيوصّي إنهم يدخّلوا \VUgras{المشط} \textsuperscript{(40)} و\VUgras{الفاس}
\textsuperscript{(39)} تحت السقف، اللي نسيوهم جنب \VUgras{عشّة} \textsuperscript{(38)}
\VUgras{الكلب} \textsuperscript{(37)}.

%%K t07-c2-09-1 p
9. --- وصوته بيخضّ \VUgras{حمامة} \textsuperscript{(41)}، اللي بتطير بكل جناحها
لجوّه \VUgras{برج الحمام} \textsuperscript{(42)}، ويخضّ \VUgras{الفراخ}
\textsuperscript{(43)} اللي بتستعجل ترجع \VUgras{الخُمّ} \textsuperscript{(44)}.

%%K t07-c2-10-1 p
10. --- وقدام \VUgras{الزريبة} \textsuperscript{(48)}، أهو \VUgras{الراعي}
\textsuperscript{(45)} العجوز راجع بـ\VUgras{قطيعه} \textsuperscript{(49)}. وماسك
\VUgras{عصايته} \textsuperscript{(46)} اللي ما بتفارقهوش أبدًا، زيها زي
\VUgras{بلوزته} \textsuperscript{(47)} اللي راح لونها؛ وبيسوق للحظيرة
\VUgras{الخرفان} \textsuperscript{(50)}، و\VUgras{النعاج} \textsuperscript{(53)}،
و\VUgras{الحملان} \textsuperscript{(54)}، و\VUgras{المعزات} \textsuperscript{(51)}،
و\VUgras{الجديان} \textsuperscript{(52)}. و\VUgras{كلبه} \textsuperscript{(55)} الوفي
أرجوس جاي في الآخر، وبنبحه بيستعجل اللي بيتأخّروا.

%%K t07-c2-11-1 p
11. --- وكل الدوشة والحركة دي ما بتقلقش هدوء \VUgras{البقرة الحلوب}
\textsuperscript{(56)} الطيبة، اللي \VUgras{جرسها} \textsuperscript{(57)} بيرنّ وهما
بيحلبوها قدام باب \VUgras{الحظيرة} \textsuperscript{(58)}.

%%K t07-c2-12-1 p
12. وهي كأنها فاهمة إنها لازم تدّي لبنها، علشان \VUgras{صاحبة العزبة}
\textsuperscript{(60)} تقدر تحضّر \VUgras{زبدة} في \VUgras{الخضّاضة}
\textsuperscript{(61)} أو \VUgras{الجبنة} \textsuperscript{(64)} الممتازة، اللي بتنقّط
من اللوح المحطوط على \VUgras{الحوض الكبير} \textsuperscript{(63)}.

%%K t07-c2-13-1 p
13. --- و\VUgras{محل اللبن} \textsuperscript{(59)} كله بيفضل نضيف أوي بفضل
\VUgras{عامل العزبة} \textsuperscript{(65)} اللي حالًا غسل \VUgras{أواني اللبن}
\textsuperscript{(62)} في الجردل الكبير اللي شايله في إيده.

%%K t07-tit-9 sub
\VUsaut{3.0mm}
\VUpk{10.2pt}{40}{الخُمّ.}
\VUsaut{3.0mm}

%%K t07-c2-14-1 p
14. --- وبقباقيبه التخينة، بيمشي بتقل شوية، وكده بيخضّ \VUgras{الديك الرومي}
\textsuperscript{(68)}، و\VUgras{البطّات} \textsuperscript{(66)}، و\VUgras{البطّ}
\textsuperscript{(67)}، و\VUgras{الوز} \textsuperscript{(69)} اللي بيفتحوا
\VUgras{مناقيرهم} \textsuperscript{(70)} على الآخر ويصوّتوا بقلق.

%%K t07-c2-14-2 p
و\VUgras{الديك} \textsuperscript{(71)}، اللي بيحبّ الخناقة أكتر، بيرفع \VUgras{عرفه}
\textsuperscript{(72)} الأحمر، وينفض ريش \VUgras{ديله} \textsuperscript{(73)}، ويصوّت
«كوكوريكو» بصوت عالي أوي.

%%K t07-c2-15-1 p
15. --- و\VUgras{الفرخة الأم} \textsuperscript{(74)} بتنقر على \VUgras{الزبل}
\textsuperscript{(81)} مع \VUgras{كتاكيتها} \textsuperscript{(75)}.
و\VUgras{الخنزير الصغير} \textsuperscript{(78)} بيجري لأمه، \VUgras{الخنزيرة}
\textsuperscript{(77)} الضخمة اللي شايفينها جنب زريبة الخنازير.

%%K t07-c2-16-1 p
16. --- وفي أقفاصهم، \VUgras{الأرانب} \textsuperscript{(79)} بتاكل بشراهة
\VUgras{الحشيش} \textsuperscript{(80)} الطري اللي بتجيبه لهم بنت صاحب العزبة.
ويوانّا بتحبّ أرانبها أوي، وكل يوم، بعد ما تلمّ \VUgras{البيض} \textsuperscript{(76)}
الطازة من الخُمّ، بتيجي تزور أصحابها الصغيّرين.

\VUsaut{6.0mm}
\VUfilet{22mm}
